% Source: https://artofproblemsolving.com/community/c3h2934037_2015_emcc_guts_round_p13_p24__exeter_math_club_competition
% Stable identity: aops-topic-2934037-q13
Horizontal translation reduces the root sum by $2B$, while vertical translation does not change the axis of symmetry. Hence $A=1028-2B$. The reflected rays $OX,OY$ differ by $90^\circ$, so $B=2C^2$. Since $C$ is a counting answer, $A,B,C$ are integers, and the interior-angle condition gives $0<A<45$. Therefore $B=\dfrac{1028-A}{2}>A$. The circle of radius $A$ about the second point meets circles of integer radii about the first exactly for radii $B-A,\ldots,B+A$. The two endpoints each give one intersection, and the $2A-1$ interior radii each give two. Thus $C=4A$. Substitution gives $64A^2+A-1028=0$, or $(A-4)(64A+257)=0$. Consequently $\boxed{(A,B,C)=(4,512,16)}$. These values satisfy every condition; a quadratic with a sufficiently negative minimum ensures both root pairs are real.
\begin{asy}
usepackage("amsmath");
size(7cm);
pair O=(0,0),P=2*dir(20),X=2*dir(-20),Y=2*dir(70);draw(O--(3,0)^^O--3*dir(45));draw(O--X--Y--cycle,dashed);draw(O--P);label("$O$",O,SW);label("$P$",P,E);label("$X$",X,S);label("$Y$",Y,N);label("schematic reflections",(1,-1.2));
\end{asy}

\begin{asy}
usepackage("amsmath");
size(7cm);
pair X=(0,0),Y=(4,0);draw(circle(Y,1));draw(circle(X,3),dashed);draw(circle(X,4),blue);draw(circle(X,5),dashed);draw(X--Y);label("$X$",X,SW);label("$Y$",Y,SE);label("endpoint and interior radii",(2,-5.5));
\end{asy}
